A comet of mass $\alpha m$ is heading (``falls'') radially towards the Sun. It
is known that the total mechanical energy of the comet is zero. The comet
crashes into Venus, whose mass is $m$. We further assume that the orbit of
Venus, before the collision, is circular with radius $R_0$. After the crash,
the comet and Venus form a single object, called ``Venus-2''.

\ioaafig{2021_TH_Q14-f1.pdf}

\begin{description}
\item[(14.1)] Find the expression in terms of $M_{sun}$ and $R_0$ for the
  orbital speed, $v_0$, of Venus before the collision. \hfill[1.0\,pt]
\item[(14.2)] Find an expression for the total mechanical energy of Venus in
  its orbit before colliding with the comet. \hfill[1.0\,pt]
\item[(14.3)] Find an expression for the radial velocity, $v_r$, the angular
  momentum, $L$, of ``Venus-2'' immediately after the collision.
  \hfill[10.0\,pt]
\item[(14.4)] Find an expression for the mechanical energy of the combined
  object ``Venus-2'' and express it in terms of energy before the collision,
  $E_i$, and $\alpha$. \hfill[5.0\,pt]
\item[(14.5)] Show that the post-collision orbit of ``Venus-2'' is elliptical
  and determine the semi-major axis of the orbit. \hfill[5.0\,pt]
\item[(14.6)] Determine if the year for the inhabitants of ``Venus-2'' has
  been shortened or lengthened because of collision with the comet. Write the
  ratio between the period of Venus-2 and Venus. \hfill[3.0\,pt]
\item[(14.7)] What should be the value of $\alpha$ such that the
  post-collision orbit of Venus-2 would make it crash in the Sun? We will call
  this as $\alpha_c$ \hfill[5.0\,pt]
\item[(14.8)] A comet with $\alpha = \alpha_c$ collided with Venus. Calculate
  the percentage change in the magnitude of Venus' velocity ($\delta v$) and
  the change in the direction of the velocity vector ($\delta\theta$)
  immediately after the collision. \hfill[5.0\,pt]
\end{description}
